\documentclass[11pt,a4paper]{article}

% ============================================================================
% PACKAGES
% ============================================================================
\usepackage[utf8]{inputenc}
\usepackage[T1]{fontenc}
\usepackage{amsmath,amssymb}
\usepackage{graphicx}
\usepackage{booktabs}
\usepackage{geometry}
\usepackage{hyperref}
\usepackage{xcolor}
\usepackage{float}

\geometry{margin=1in}

\begin{document}

% ============================================================================
% METHODS SECTION
% ============================================================================
\section*{Methods}

\subsection*{Study Design and Data Source}

We conducted a retrospective analysis of drug-drug interactions (DDIs) extracted from DrugBank version 5.1.9. The study focused on cardiovascular and antithrombotic therapeutic categories, comprising 759,774 unique DDI pairs involving 1,247 drugs (892 cardiovascular, 355 antithrombotic agents). Interaction descriptions and initial severity predictions were obtained using zero-shot classification.

\subsection*{Baseline Severity Classification}

Initial severity predictions were generated using the BART-MNLI model (\texttt{facebook/bart-large-mnli}) via zero-shot natural language inference. The model classified each DDI into four severity categories: Contraindicated, Major, Moderate, or Minor, based on the interaction description text.

\subsection*{Hybrid Recalibration Framework}

To address systematic over-classification bias observed in zero-shot predictions, we developed a hybrid evidence-based severity recalibration (HEBSR) framework. The framework combines three complementary scoring mechanisms through a weighted ensemble:

\begin{equation}
S_{\text{final}} = w_s \cdot S_{\text{semantic}} + w_c \cdot S_{\text{confidence}} + w_d \cdot S_{\text{drug\_class}}
\label{eq:hybrid}
\end{equation}

where $w_s = 0.45$, $w_c = 0.25$, and $w_d = 0.30$. The semantic component receives highest weight (0.45) because embedding-based similarity generalizes across diverse DDI descriptions without relying on fixed keyword patterns. Confidence adjustment receives lowest weight (0.25) as a corrective factor applying only to low-confidence predictions. Drug class scoring receives intermediate weight (0.30) to provide pharmacological priors that complement description-level analysis. Weights were determined via grid search (0.1--0.6, step 0.05) minimizing Jensen-Shannon divergence from DDInter targets while maintaining $>$95\% sensitivity for Black Box Warning pairs.

\subsubsection*{Component 1: Semantic Similarity Scoring}

We employed sentence embeddings to compare interaction descriptions against severity-specific prototype phrases. Using the \texttt{all-MiniLM-L6-v2} transformer model, each interaction description was encoded into a 384-dimensional vector. Severity was assigned based on cosine similarity to curated prototype centroids representing each severity class:

\begin{equation}
\text{sim}(d, k) = \frac{\mathbf{e}_d \cdot \mathbf{c}_k}{\|\mathbf{e}_d\| \|\mathbf{c}_k\|}
\end{equation}

where $\mathbf{e}_d$ is the embedding of description $d$ and $\mathbf{c}_k$ is the centroid for severity class $k$. Severity prototypes were designed to capture canonical clinical scenarios for each category, including life-threatening outcomes (Contraindicated), serious adverse events requiring intervention (Major), pharmacokinetic alterations requiring monitoring (Moderate), and clinically insignificant effects (Minor).

\subsubsection*{Component 2: Confidence-Weighted Adjustment}

The original zero-shot prediction confidence was incorporated to penalize low-confidence high-severity predictions:

\begin{equation}
S_{\text{confidence}} = 
\begin{cases}
3.0 & \text{if } \hat{y} = \texttt{Contraindicated} \land c < 0.65 \\
2.5 & \text{if } \hat{y} \in \{\texttt{Contraindicated, Major}\} \land c < 0.50 \\
\phi(\hat{y}) & \text{otherwise}
\end{cases}
\end{equation}

where $\hat{y}$ is the predicted label, $c$ is prediction confidence, and $\phi$ maps labels to numeric scores (Contraindicated=4, Major=3, Moderate=2, Minor=1).

\subsubsection*{Component 3: Drug Class Risk Profiling}

Pharmacological class membership informed severity adjustment based on known high-risk drug combinations. High-risk classes included anticoagulants, antiplatelets, QT-prolonging agents, MAOIs, and narrow therapeutic index drugs. Interactions involving multiple high-risk classes received elevated severity scores.

\subsection*{Severity Threshold Calibration}

The composite score $S_{\text{final}}$ was mapped to severity categories using thresholds calibrated to match DDInter-derived target distributions:

\begin{equation}
\text{Severity} = 
\begin{cases}
\texttt{Contraindicated} & S_{\text{final}} \geq 3.2 \\
\texttt{Major} & 2.5 \leq S_{\text{final}} < 3.2 \\
\texttt{Moderate} & 2.0 \leq S_{\text{final}} < 2.5 \\
\texttt{Minor} & S_{\text{final}} < 2.0
\end{cases}
\end{equation}

Target distributions were derived from clinical literature: Contraindicated (5\%), Major (25\%), Moderate (60\%), Minor (10\%).

\subsection*{Known Pair Override}

A curated set of 34 DDI pairs with established FDA warnings or clinical guideline contraindications bypassed the hybrid scoring mechanism and were assigned their clinically-validated severity directly.

\subsection*{Implementation}

The recalibration pipeline was implemented in Python using PyTorch for GPU acceleration and the SentenceTransformers library for embedding computation. Processing was performed on an NVIDIA RTX PRO 5000 GPU (48GB VRAM) with 24 CPU cores.

\subsection*{Validation}

Clinical validity was assessed through: (1) high-risk combination sensitivity---the proportion of known dangerous combinations classified as Major or Contraindicated; (2) correlation with TWOSIDES Proportional Reporting Ratio scores from real-world adverse event reports; and (3) Cohen's kappa agreement with clinical pharmacist assessments on a validation subset.

\subsection*{Statistical Analysis}

Distribution alignment was quantified using Jensen-Shannon divergence between predicted and target distributions. Chi-square goodness-of-fit tests assessed statistical deviation from targets. All analyses were conducted using SciPy v1.11.

\end{document}
